\documentclass[11pt]{article}
\usepackage[margin=1in]{geometry}
\usepackage{amsmath,amssymb,booktabs,hyperref}
\title{PINNeAPPle: Towards Reproducible and Physics-Verified Physics-Informed Learning}
\author{Yan Barros \\and Enrico Schiassi \\\\ \\small Affiliations to be confirmed}
\date{Draft -- not for submission}
\begin{document}
\maketitle
\begin{abstract}
Physics-informed learning frameworks vary in model coverage, experiment
reproducibility, and treatment of physical validity. We describe a protocol
for evaluating PINNeAPPle against established frameworks using matched
partial differential equation benchmarks, operator-learning tasks, and
independent physical checks. Quantitative conclusions will be included
only after the preregistered experiments have been executed.
\end{abstract}
\section{Introduction}
Reproducible physical machine learning requires more than predictive
accuracy: governing equations, boundary conditions, provenance, and
validity limits must be explicit. This paper studies the design and
empirical evaluation of PINNeAPPle.
\section{Related Work}
Discuss DeepXDE, NVIDIA PhysicsNeMo, neuraloperator, PINNs, DeepONet,
Fourier neural operators and scientific ML verification, with verified
primary citations before submission.
\section{Software Architecture}
Document only features verified in the pinned release, with code paths.
\section{Research Questions and Experimental Protocol}
See the accompanying protocol. Separate matched-model and native-framework
comparisons. Report independent test residuals and domain-shift experiments.
\section{Results}
\textbf{Pending execution.} No numerical results are asserted in this draft.
\section{Threats to Validity and Limitations}
Include benchmark selection bias, numerical-reference uncertainty,
compute budget, and failure cases.
\section{Reproducibility and Artifact Availability}
Provide immutable source revision, data checksums, environments,
scripts, and archival DOI prior to submission.
\section{Conclusion}
To be written after experimental evaluation.
\end{document}
